% AA-16 | Algebra abstrakcyjna | audyt II
\setlength{\parskip}{3pt}
\chapter{Moduły, homomorfizmy i ciągi dokładne}
\label{aa:ch:modules-introduction}

\section{Dlaczego po grupach i pierścieniach pojawiają się moduły?}
Przestrzeń wektorowa nad ciałem $K$ jest grupą abelową,
na której elementy $K$ działają jak skalary.
Zastępując ciało pierścieniem z jedynką, otrzymujemy
pojęcie modułu. Ta zmiana nie jest kosmetyczna: nawet
najprostsze wyniki teorii przestrzeni wektorowych,
takie jak istnienie bazy dla każdej przestrzeni, mogą
przestać obowiązywać. Z drugiej strony jeden język
obejmie grupy abelowe, przestrzenie wektorowe,
ideały pierścieni, przedstawienia grup oraz układy
równań liniowych o współczynnikach z pierścienia.

W tym rozdziale $R$ jest pierścieniem z jedynką,
niekoniecznie przemiennym. W definicjach lewego
modułu zapisujemy działanie skalara z lewej strony.

\begin{definicja}[Lewy i prawy moduł nad pierścieniem]
\label{aa:def:modules-full}
\emph{Lewym $R$-modułem} nazywamy grupę abelową
$(M,+,0_M)$ wraz z działaniem $R\times M\to M$,
$(r,m)\mapsto rm$, takim, że dla wszystkich
$r,s\in R$ i $m,n\in M$:
\begin{enumerate}[label=\textup{(M\arabic*)}]
\item $(r+s)m=rm+sm$;
\item $r(m+n)=rm+rn$;
\item $(rs)m=r(sm)$;
\item $1_Rm=m$.
\end{enumerate}
\emph{Prawym $R$-modułem} nazywamy grupę abelową $M$
z działaniem $M\times R\to M$, $(m,r)\mapsto mr$,
dla którego, dla dowolnych $m,n\in M$ i $r,s\in R$,
zachodzą
$(m+n)r=mr+nr$, $m(r+s)=mr+ms$,
$(mr)s=m(rs)$ oraz $m1_R=m$.
Jeśli $R$ jest przemienny, z każdego lewego
$R$-modułu otrzymujemy prawy, kładąc $mr=rm$:
warunek $(mr)s=m(rs)$ wynika wtedy z
$(sr)m=s(rm)$ i przemienności skalarów.
\end{definicja}

\begin{stwierdzenie}[Pierwsze konsekwencje aksjomatów]
\label{aa:st:modules-zero-laws}
Niech $M$ będzie lewym $R$-modułem. Dla każdego
$r\in R$ i $m\in M$ zachodzą
$0_Rm=0_M$, $r0_M=0_M$,
$(-r)m=-(rm)$ oraz $r(-m)=-(rm)$.
\end{stwierdzenie}
\begin{proof}
Z $(0_R+0_R)m=0_Rm+0_Rm$ otrzymujemy
$0_Rm=0_M$ po odjęciu $0_Rm$ w grupie addytywnej
$M$. Analogicznie
$r(0_M+0_M)=r0_M+r0_M$ daje $r0_M=0_M$.
Ponieważ $0_Rm=(r+(-r))m=rm+(-r)m$, element
$(-r)m$ jest przeciwny do $rm$. Ostatnia równość
wynika tak samo z $r(m+(-m))=r0_M=0_M$.
\end{proof}

\begin{przyklad}[Cztery zasadniczo różne moduły]
\label{aa:prz:modules-four}
\begin{enumerate}[label=\textup{\arabic*)}]
\item Każda grupa abelowa $A$ jest kanonicznie
$\mathbb Z$-modułem, ponieważ $na$ oznacza
$n$-krotną sumę elementu $a$. Homomorfizmy takich
modułów są dokładnie homomorfizmami grup abelowych.
\item Każda przestrzeń wektorowa nad ciałem $K$
jest $K$-modułem. Dla $K$ wszystkie moduły mają bazy,
co udowadnia standardowe twierdzenie o bazie
przestrzeni wektorowej.
\item Każdy ideał lewostronny $I\subseteq R$ jest
lewym $R$-modułem z mnożeniem odziedziczonym z $R$.
Przykładowo $2\mathbb Z$ jest $\mathbb Z$-modułem
izomorficznym z $\mathbb Z$, mimo że jako podzbiór
nie zawiera elementu $1$.
\item Jeśli $R=\mathbb Z/4\mathbb Z$, to
$M=2R=\{\bar0,\bar2\}$ jest $R$-modułem.
Nie jest wolnym $R$-modułem: dla dowolnego
niezerowego $m\in M$ zachodzi $\bar2m=0$,
więc żaden niezerowy element $M$ nie może stanowić
jednoelementowej bazy. Co więcej, każdy niezerowy
moduł wolny nad $R$ zawiera kopię $R$: jeśli $e$
należy do jego bazy, odwzorowanie $r\mapsto re$ jest
różnowartościowe dzięki jednoznaczności współrzędnych.
Taki moduł ma zatem co najmniej cztery elementy,
podczas gdy $M$ ma tylko dwa.
\end{enumerate}
\end{przyklad}

\begin{definicja}[Podmoduł i homomorfizm modułów]
\label{aa:def:modules-sub-hom}
Niech $M,N$ będą lewymi $R$-modułami.
\emph{Podmoduł} $L\subseteq M$ to podgrupa addytywna
zamknięta na działanie każdego $r\in R$.
\emph{Homomorfizm $R$-modułów} to odwzorowanie
$f:M\to N$ spełniające $f(m+m')=f(m)+f(m')$
i $f(rm)=rf(m)$ dla wszystkich $m,m'\in M$, $r\in R$.
Zbiór wszystkich takich homomorfizmów oznaczamy
$\operatorname{Hom}_R(M,N)$. W szczególności każdy homomorfizm
jest addytywny, więc $f(0_M)=0_N$ oraz $f(-m)=-f(m)$:
pierwsza równość wynika z $f(0_M)=f(0_M)+f(0_M)$,
a druga z $f(m)+f(-m)=f(0_M)$.
\end{definicja}

\begin{stwierdzenie}[Grupa homomorfizmów modułów]
\label{aa:st:modules-hom-abelian}
Dla lewych $R$-modułów $M,N$ zbiór
$\operatorname{Hom}_R(M,N)$ jest grupą abelową przy
punktowym dodawaniu $(f+g)(m)=f(m)+g(m)$.
Jeśli $R$ jest przemienny, staje się dodatkowo
$R$-modułem dla $(rf)(m)=r f(m)$.
\end{stwierdzenie}
\begin{proof}
Dla homomorfizmów $f,g$ ich suma zachowuje dodawanie,
ponieważ dla $m,m'\in M$ mamy
$(f+g)(m+m')=f(m)+f(m')+g(m)+g(m')$
$=(f+g)(m)+(f+g)(m')$; środkowa zamiana
składników jest dopuszczalna, gdyż $N$ jest grupą abelową.
Dla $r\in R$ oraz $m\in M$ mamy
$(f+g)(rm)=f(rm)+g(rm)=rf(m)+rg(m)
=r(f(m)+g(m))$, zatem $f+g$ jest $R$-liniowe.
Odwzorowanie zerowe i odwzorowanie $-f$ są także
$R$-liniowe. Łączność i przemienność dodawania
pochodzą z grupy addytywnej $N$.
Jeśli $R$ jest przemienny, to
$(af)(rm)=a f(rm)=a(rf(m))=r(af(m))$;
pozostałe aksjomaty modułu wynikają z aksjomatów $N$.
Przemienność jest tu istotna: dla nieprzemiennego
$R$ wzór $(af)(m)=af(m)$ może nie zachowywać
$R$-liniowości, ponieważ $ar$ nie musi być równe $ra$.
\end{proof}

\begin{lemat}[Kryterium podmodułu]
\label{aa:lem:modules-submodule-test}
Niepusty podzbiór $L\subseteq M$ jest podmodułem
lewego $R$-modułu $M$ wtedy i tylko wtedy, gdy
$\ell+\ell'\in L$ i $r\ell\in L$ dla
wszystkich $\ell,\ell'\in L$ i $r\in R$.
\end{lemat}
\begin{proof}
Konieczność wynika z definicji. Przy założonej
zamkniętości wybierzmy $\ell\in L$.
Ponieważ $0_R\ell=0_M$, do $L$ należy zero.
Dalej $-\ell=(-1_R)\ell$ należy do $L$.
Zamkniętość na dodawanie i elementy przeciwne
pokazuje, że $L$ jest podgrupą addytywną;
zamkniętość na działanie $R$ już założyliśmy.
\end{proof}

\begin{twierdzenie}[Jądro, obraz i pierwszy izomorfizm modułów]
\label{aa:tw:modules-first-isomorphism}
Dla homomorfizmu lewych $R$-modułów $f:M\to N$
zbiór $\ker f$ jest podmodułem $M$, a
$\operatorname{im}f$ podmodułem $N$. Wzór
$\bar f(m+\ker f)=f(m)$ wyznacza izomorfizm
$R$-modułów $M/\ker f\cong\operatorname{im}f$.
\end{twierdzenie}
\begin{proof}
Jeżeli $m,m'\in\ker f$ i $r\in R$, to
$f(m+m')=f(m)+f(m')=0$ oraz $f(rm)=rf(m)=0$.
Jądro zawiera $0_M$, gdyż $f(0_M)=0_N$.
Jeśli $n=f(m)$ i $n'=f(m')$ należą do obrazu,
to $n+n'=f(m+m')$ i $rn=f(rm)$ także do niego należą.
Obraz jest niepusty, ponieważ zawiera $f(0_M)=0_N$.
Lemat~\ref{aa:lem:modules-submodule-test} pozwala zatem
zastosować kryterium podmodułu zarówno do jądra,
jak i do obrazu.

Grupa addytywna $M/\ker f$ jest ilorazem przez
podgrupę; definiujemy działanie $r(m+\ker f)=rm+\ker f$.
Jeśli $m'=m+k$ z $k\in\ker f$, to
$rm'=rm+rk$ oraz $rk\in\ker f$, więc działanie
nie zależy od reprezentanta. Aksjomaty modułu
przechodzą z $M$ na iloraz. Ponieważ
$f(m)=f(m')$ wtedy i tylko wtedy, gdy
$m-m'\in\ker f$, odwzorowanie $\bar f$
jest dobrze określone i różnowartościowe.
Jest na obraz z definicji. Zachowuje dodawanie
i skalary, ponieważ $f$ je zachowuje.
\end{proof}

\begin{twierdzenie}[Drugi i trzeci izomorfizm modułów]
\label{aa:tw:modules-second-third}
Niech $M$ będzie lewym $R$-modułem.
\begin{enumerate}[label=\textup{(\roman*)}]
\item Dla podmodułów $A,B\subseteq M$ zachodzi
$(A+B)/B\cong A/(A\cap B)$.
\item Jeśli $A\subseteq B\subseteq M$ są podmodułami,
to $(M/A)/(B/A)\cong M/B$.
\end{enumerate}
\end{twierdzenie}
\begin{proof}
W (i) suma $A+B=\{a+b:a\in A,\ b\in B\}$
jest podmodułem $M$: suma dwóch takich elementów,
ich elementy przeciwne i wszystkie wielokrotności
przez $R$ znów należą do $A+B$.
Rozważmy homomorfizm $A\to(A+B)/B$,
$a\mapsto a+B$. Jest on na, ponieważ każda klasa
$a+b+B$ jest równa $a+B$. Jego jądro to dokładnie
$A\cap B$. Z pierwszego twierdzenia o izomorfizmie
otrzymujemy $(i)$. W (ii) homomorfizm
$M/A\to M/B$ dany przez $m+A\mapsto m+B$
jest dobrze określony, bo $A\subseteq B$, i jest
surjektywny. Jego jądro składa się dokładnie z
klas $b+A$, $b\in B$, czyli jest równe $B/A$.
Pierwsze twierdzenie o izomorfizmie daje (ii).
\end{proof}

\section{Moduły wolne i prezentacje}
\begin{definicja}[Moduł wolny]
\label{aa:def:modules-free}
Rodzina $(e_s)_{s\in S}$ elementów lewego $R$-modułu
$F$ jest \emph{bazą}, jeśli każdy $x\in F$ ma dokładnie
jedno przedstawienie $x=\sum_{s\in S}r_se_s$,
w którym tylko skończenie wiele współczynników
$r_s\in R$ jest niezerowych. Moduł mający bazę
nazywamy \emph{wolnym}. Symbol $R^{(S)}$ oznacza
zbiór funkcji $S\to R$ o skończonym nośniku,
z dodawaniem i lewym działaniem $R$ określonymi
współrzędnie. Dla $s\in S$ niech $e_s$ przyjmuje
wartość $1_R$ w punkcie $s$ i $0_R$ poza nim.
Wówczas każdy $x\in R^{(S)}$ ma jednoznaczny zapis
$x=\sum_{s\in S}x(s)e_s$; równość współrzędnych
w punkcie $s$ dowodzi jednoznaczności. W szczególności
$(e_s)_{s\in S}$ jest bazą, a dla $S=\varnothing$
otrzymujemy moduł zerowy.
\end{definicja}

\begin{twierdzenie}[Własność uniwersalna modułu wolnego]
\label{aa:tw:modules-free-universal}
Niech $S$ będzie zbiorem, a $M$ lewym $R$-modułem.
Dla każdej funkcji $u:S\to M$ istnieje dokładnie
jeden homomorfizm $R$-modułów
$\bar u:R^{(S)}\to M$, który wysyła standardowy
wektor bazowy $e_s$ na $u(s)$ dla każdego $s\in S$.
\end{twierdzenie}
\begin{proof}
Dla funkcji $x:S\to R$ o skończonym nośniku
określamy $\bar u(x)=\sum_{s\in S}x(s)u(s)$.
Suma zawiera tylko skończenie wiele niezerowych
składników, więc jest dobrze określona.
Jej addytywność wynika z rozdzielności działania,
a $\bar u(rx)=r\bar u(x)$ z aksjomatu
$(rs)m=r(sm)$: istotna jest tu kolejność
mnożenia skalarów w nieprzemiennym $R$.
Współrzędne $e_s$ dają $\bar u(e_s)=u(s)$.
Jeżeli $v$ jest innym takim homomorfizmem,
to dla $x=\sum_s x(s)e_s$ mamy
$v(x)=\sum_s x(s)v(e_s)=\bar u(x)$.
Stąd jednoznaczność.
\end{proof}

\begin{definicja}[Prezentacja modułu przez generatory i relacje]
\label{aa:def:modules-presentation}
\emph{Prezentacją} lewego $R$-modułu $M$ nazywamy
ciąg homomorfizmów
$R^{(T)}\xrightarrow{d}R^{(S)}\xrightarrow{p}M\to0$
taki, że $p$ jest surjekcją i $\operatorname{im}d=\ker p$.
Oznacza to, że elementy $p(e_s)$ generują $M$,
a elementy $d(e_t)$ są relacjami generującymi
wszystkie zależności między generatorami.
\end{definicja}

\begin{twierdzenie}[Każdy moduł ma prezentację wolną]
\label{aa:tw:modules-presentation-exists}
Dla każdego lewego $R$-modułu $M$ istnieją zbiory
$S,T$ i prezentacja $R^{(T)}\to R^{(S)}\to M\to0$.
\end{twierdzenie}
\begin{proof}
Weźmy $S=M$ jako zbiór i funkcję identycznościową
$u:S\to M$. Własność uniwersalna daje homomorfizm
$p:R^{(S)}\to M$ wysyłający $e_m$ na $m$.
Jest on surjektywny, bo każdy $m\in M$ jest obrazem
$e_m$. Niech $K=\ker p$ i weźmy $T=K$ jako zbiór.
Własność uniwersalna modułu wolnego daje
$d:R^{(T)}\to R^{(S)}$ określony na standardowych
wektorach przez $d(e_k)=k$. Wtedy
$\operatorname{im}d=K$; istotnie, obrazy należą do $K$,
a każdy $k\in K$ jest obrazem $e_k$.
Otrzymujemy wymaganą prezentację.
\end{proof}

\begin{przyklad}[Prezentacje modułów nad liczbami całkowitymi]
\label{aa:prz:modules-presentations-z}
Dla $n\ge2$ grupa cykliczna $\mathbb Z/n\mathbb Z$
ma prezentację
$\mathbb Z\xrightarrow{\cdot n}\mathbb Z\to\mathbb Z/n\mathbb Z\to0$.
Generator to klasa liczby $1$, a jedyna relacja
ma postać $n[1]=0$. Dla
$\mathbb Z/6\mathbb Z\oplus\mathbb Z/10\mathbb Z$
możemy użyć dwóch generatorów $u,v$ i relacji
$6u=0$, $10v=0$: homomorfizm
$\mathbb Z^2\to\mathbb Z^2$ o macierzy
$\operatorname{diag}(6,10)$ opisuje wszystkie relacje:
jej obraz wynosi $6\mathbb Z\oplus10\mathbb Z$,
czyli jądro epimorfizmu
$(a,b)\mapsto(a+6\mathbb Z,b+10\mathbb Z)$.
Gdy przejdziemy do postaci Smitha, znajdziemy
kanoniczny sposób upraszczania takich prezentacji.
\end{przyklad}

\begin{przyklad}[Moduł z jednym generatorem nie musi być wolny]
\label{aa:prz:modules-cyclic-not-free}
Niech $K$ będzie ciałem, a $R=K[t]/(t^2)$.
Moduł $M=R/(t)$ jest generowany przez klasę $1$,
ale $t\ne0$ w $R$ i $t\cdot(1+(t))=0$ w $M$.
W pierścieniu $R$ klasa $t$ jest niezerowa,
gdyż $t\notin(t^2)\subseteq K[t]$. Zeruje ona cały
$M$, ponieważ $t\cdot(r+(t))=tr+(t)=0$
dla dowolnego $r\in R$ (ideał $(t)$ jest dwustronny).
Gdyby $M$ był wolnym $R$-modułem z niepustą bazą,
to żaden niezerowy skalar nie mógłby zerować
wszystkich elementów: dla elementu bazowego $e_s$
równość $te_s=0$ przeczyłaby jednoznaczności
współrzędnych. Zatem moduł nie jest wolny. Pokazuje to, że generator
i element bazy to różne pojęcia.
\end{przyklad}

\section{Suma prosta, iloczyn prosty i ciągi dokładne}
\begin{definicja}[Suma prosta i iloczyn prosty modułów]
\label{aa:def:modules-direct-sum}
Dla rodziny lewych $R$-modułów $(M_i)_{i\in I}$
\emph{iloczyn prosty} $\prod_{i\in I}M_i$ składa się
ze wszystkich rodzin elementów, z działaniem
$R$ i dodawaniem określonymi współrzędnie.
\emph{Suma prosta} $\bigoplus_{i\in I}M_i$ jest
podmodułem tych rodzin, które mają skończony nośnik.
Istotnie, nośnik sumy dwóch rodzin jest zawarty
w sumie mnogościowej ich nośników, a nośnik
rodziny pomnożonej przez skalar jest zawarty
w nośniku wyjściowej rodziny. Oba pozostają więc
skończone. Rodzina zerowa należy do sumy prostej,
a element przeciwny nie zwiększa nośnika.
Dla skończonego zbioru indeksów obie konstrukcje
są równe, ale dla nieskończonego mogą się różnić.
Na przykład dla $I=\mathbb N$ i $M_i=\mathbb Z$
rodzina $(1,1,\ldots)$ należy do
$\prod_{i\in\mathbb N}\mathbb Z$, ale nie do
$\bigoplus_{i\in\mathbb N}\mathbb Z$, ponieważ
ma nieskończony nośnik.
\end{definicja}

\begin{definicja}[Dokładność ciągu modułów]
\label{aa:def:modules-exact}
Ciąg $L\xrightarrow{f}M\xrightarrow{g}N$
homomorfizmów lewych $R$-modułów jest \emph{dokładny
w $M$}, jeśli $\operatorname{im}f=\ker g$.
\emph{Krótki ciąg dokładny} ma postać
\[
0\longrightarrow A\xrightarrow{i}B
 \xrightarrow{p}C\longrightarrow0,
\]
gdzie $i$ jest iniekcją, $p$ jest surjekcją oraz
$\operatorname{im}i=\ker p$.
\end{definicja}

\begin{twierdzenie}[Kryteria rozszczepienia ciągu dokładnego]
\label{aa:tw:modules-splitting}
Dla krótkiego ciągu dokładnego lewych $R$-modułów
$0\to A\xrightarrow{i}B\xrightarrow{p}C\to0$
następujące warunki są równoważne:
\begin{enumerate}[label=\textup{(\roman*)}]
\item istnieje $R$-liniowe $s:C\to B$ takie, że
$p\circ s=\operatorname{id}_C$;
\item istnieje $R$-liniowe $r:B\to A$ takie, że
$r\circ i=\operatorname{id}_A$;
\item istnieje izomorfizm $B\cong A\oplus C$
zgodny z danym ciągiem, to znaczy przeprowadzający
$i(a)$ na $(a,0)$ i $p$ na rzut na drugą współrzędną.
\end{enumerate}
\end{twierdzenie}
\begin{proof}
Załóżmy (i). Dla $b\in B$ element
$b-s(p(b))$ należy do $\ker p=\operatorname{im}i$.
Ponieważ $i$ jest iniekcją, istnieje dokładnie
jedno $a\in A$ takie, że
$i(a)=b-s(p(b))$. Wzór
$\Phi(a,c)=i(a)+s(c)$ określa $R$-liniowe
odwzorowanie $A\oplus C\to B$. Jest ono na:
dla $b\in B$ wybieramy wyżej opisane $a$ i
$c=p(b)$. Jeśli $\Phi(a,c)=0$, po zastosowaniu
$p$ dostajemy $c=0$, a następnie $i(a)=0$,
więc $a=0$. Zatem $\Phi$ jest izomorfizmem.
Jego odwrotność ma postać
$\Phi^{-1}(b)=(i^{-1}(b-s(p(b))),p(b))$,
gdzie $i^{-1}$ jest odwrotnością $i:A\to\operatorname{im}i$;
argument pierwszej współrzędnej należy do
$\operatorname{im}i$ na mocy dokładności ciągu.
Ponadto $\Phi(a,0)=i(a)$ i $p\Phi(a,c)=c$.
Dostaliśmy (iii).

Z (iii) otrzymujemy (ii): jeśli
$\Psi:B\to A\oplus C$ jest izomorfizmem zgodnym
z ciągiem, to definiujemy $r=\operatorname{pr}_A\circ\Psi$.
Wtedy $r(i(a))=\operatorname{pr}_A(a,0)=a$.

Załóżmy teraz (ii). Dla $b\in B$ element
$b-i(r(b))$ należy do $\ker r$. Dlatego
każdy $b\in B$ zapisuje się jako
$b=i(r(b))+(b-i(r(b)))$, przy czym drugi składnik
leży w $\ker r$. Rozkład jest jednoznaczny:
$\ker r\cap i(A)=0$, gdyż dla $i(a)\in\ker r$
mamy $0=r(i(a))=a$. Zatem
$B=i(A)\oplus\ker r$. Ograniczenie
$p|_{\ker r}:\ker r\to C$ jest iniektywne,
ponieważ jego jądro to
$\ker r\cap\ker p=\ker r\cap\operatorname{im}i=0$.
Jest też na: dla $c\in C$ wybieramy $b\in B$
z $p(b)=c$ i mamy
$p(b-i(r(b)))=c$. Niech
$s$ będzie odwrotnością tego ograniczenia,
po której stosujemy włączenie $\ker r\hookrightarrow B$.
Odwrotność jest $R$-liniowa: jeśli $c_1=p(u_1)$,
$c_2=p(u_2)$ dla $u_1,u_2\in\ker r$, to
$p(u_1+u_2)=c_1+c_2$ oraz $p(tu_1)=tc_1$
dla każdego $t\in R$. Jednoznaczność przeciwobrazów
w $\ker r$ zapewnia zachowanie dodawania i skalarów.
Wtedy $p\circ s=\operatorname{id}_C$, co daje (i).
\end{proof}

\begin{przyklad}[Krótki ciąg, który się nie rozszczepia]
\label{aa:prz:modules-nonsplit-z}
Rozważmy ciąg $\mathbb Z$-modułów
\[
0\longrightarrow\mathbb Z\xrightarrow{\cdot2}\mathbb Z
 \xrightarrow{\pi}\mathbb Z/2\mathbb Z\longrightarrow0.
\]
Jest dokładny: mnożenie przez $2$ jest iniektywne,
$\pi$ jest surjekcją, a $\ker\pi=2\mathbb Z$.
Gdyby istniał przekrój liniowy
$s:\mathbb Z/2\mathbb Z\to\mathbb Z$, to
$\pi(s([1]))=[1]$, więc $s([1])$ byłby nieparzysty.
Jednak $2s([1])=s(2[1])=s([0])=0$,
co w $\mathbb Z$ wymusza $s([1])=0$.
Sprzeczność dowodzi, że ciąg się nie rozszczepia.
W szczególności nie każdy podmoduł jest
składnikiem sumy prostej.
\end{przyklad}

\section{Moduły rzutowe}
\begin{definicja}[Moduł rzutowy]
\label{aa:def:modules-projective}
Niech $R$ będzie pierścieniem z jedynką,
niekoniecznie przemiennym. Lewy $R$-moduł $P$ jest
\emph{modułem rzutowym}, jeśli dla dowolnych lewych
$R$-modułów $E,N$, każdego surjektywnego homomorfizmu
$R$-modułów $q:E\to N$ i każdego homomorfizmu
$R$-modułów $f:P\to N$
istnieje homomorfizm $\widetilde f:P\to E$
spełniający $q\circ\widetilde f=f$.
Mówimy, że $\widetilde f$ jest \emph{podniesieniem}
$f$ względem $q$.
\end{definicja}

\begin{twierdzenie}[Charakterystyka modułów rzutowych]
\label{aa:tw:modules-projective-summand}
Niech $R$ będzie pierścieniem z jedynką,
niekoniecznie przemiennym, a $P$ lewym $R$-modułem.
Następujące warunki są równoważne:
\begin{enumerate}[label=\textup{(\roman*)}]
\item $P$ jest modułem rzutowym;
\item każdy krótki ciąg dokładny
$0\to A\to B\to P\to0$ rozszczepia się;
\item $P$ jest izomorficzny ze składnikiem sumy prostej
pewnego wolnego lewego $R$-modułu.
\end{enumerate}
\end{twierdzenie}
\begin{proof}
Załóżmy (i) i rozważmy krótki ciąg dokładny
$0\to A\to B\xrightarrow{q}P\to0$.
Stosując definicję modułu rzutowego, wybierzmy $f=\operatorname{id}_P$.
Otrzymujemy $s:P\to B$ z $q\circ s=\operatorname{id}_P$.
Twierdzenie~\ref{aa:tw:modules-splitting} daje rozszczepienie,
a więc (ii).

Załóżmy (ii). Z Twierdzenia~\ref{aa:tw:modules-presentation-exists}
istnieje surjekcja $p:F\to P$ z wolnego modułu $F$.
Krótki ciąg $0\to\ker p\to F\xrightarrow{p}P\to0$
rozszczepia się, zatem $F\cong\ker p\oplus P$.
Daje to (iii).

Załóżmy (iii). Najpierw pokażemy, że wolny moduł
$F=R^{(S)}$ jest modułem rzutowym. Dla surjekcji
$q:E\to N$ i homomorfizmu $f:F\to N$ wybierzmy
dla każdego $s\in S$ element $e_s'\in E$ taki, że
$q(e_s')=f(e_s)$. Przy nieskończonym $S$ korzystamy
tu z aksjomatu wyboru dla niepustych włókien surjekcji.
Własność uniwersalna modułu
wolnego daje $\widetilde f:F\to E$ z
$\widetilde f(e_s)=e_s'$. Wtedy $q\widetilde f$
i $f$ są równe na bazie $F$, więc są równe
na całym $F$.

Niech teraz $F=P\oplus Q$; przypadek modułu
izomorficznego z $P$ rozpatrujemy po przeniesieniu
struktury przez izomorfizm. Dla surjekcji $q:E\to N$
i homomorfizmu $f:P\to N$ rozważmy
$h:F\to N$, $h(p,q')=f(p)$. Wolność $F$
pozwala podnieść $h$ do $\widetilde h:F\to E$.
Ograniczając $\widetilde h$ do $P\oplus\{0\}$
dostajemy podniesienie $f$. Zatem $P$ jest modułem rzutowym i zachodzi (i).
\end{proof}

\begin{przyklad}[Własność bycia modułem rzutowym zależy od pierścienia]
\label{aa:prz:modules-projectivity-example}
Każdy moduł nad ciałem $K$ jest wolny,
a więc jest modułem rzutowym. Moduł $\mathbb Z/2\mathbb Z$
nie jest modułem rzutowym nad $\mathbb Z$, ponieważ
krótki ciąg dokładny z
Przykładu~\ref{aa:prz:modules-nonsplit-z} nie rozszczepia
się. Z drugiej strony dla pierścienia
$R=K\times K$ moduł $P=K\times\{0\}$ jest
rzutowy: $R=P\oplus(\{0\}\times K)$ jako
moduł nad sobą. Nie jest on wolny jako $R$-moduł:
$P\ne0$, a niezerowy element $(0,1)\in R$ zeruje
wszystkie elementy $P$. Gdyby $P$ miał niepustą bazę,
dla jej elementu $e$ mielibyśmy $(0,1)e\ne0$,
gdyż współczynnik $(0,1)$ jest niezerowy i
przedstawienie w bazie jest jednoznaczne.
Otrzymujemy sprzeczność. Jest to konkretny
przykład modułu rzutowego, który nie jest wolny.
\end{przyklad}


\begin{przyklad}[Rozszczepienie nie wynika z dokładności]
\label{aa:prz:modules-split-versus-exact}
Dla każdego pierścienia $R$ z jedynką i lewych $R$-modułów
$A,C$ ciąg
\[
0\longrightarrow A\xrightarrow{i}A\oplus C
\xrightarrow{p}C\longrightarrow0,
\qquad i(a)=(a,0),\quad p(a,c)=c,
\]
jest dokładny i rozszczepiony: odwzorowanie $s(c)=(0,c)$
spełnia $p\circ s=\operatorname{id}_C$. W przeciwieństwie
do tego ciąg z Przykładu~\ref{aa:prz:modules-nonsplit-z}
jest dokładny, ale nie ma takiego przekroju. Dokładność
wymaga jedynie równości obrazu i jądra, nie istnienia
uzupełniającego podmodułu.
\end{przyklad}

\begin{przyklad}[Dlaczego nie każdy homomorfizm można mnożyć przez skalar z lewej]
\label{aa:prz:modules-hom-noncommutative}
Niech $K$ będzie ciałem, $R=M_2(K)$, a $M=N={}_RR$.
Identyczność $f=\operatorname{id}_R$ jest lewym
$R$-liniowym endomorfizmem. Ustalmy $a=E_{12}$ i
rozważmy odwzorowanie $h(x)=af(x)=ax$. Gdyby $h$
było lewym $R$-liniowym homomorfizmem, to dla
$r=E_{21}$ i $x=1_R$ mielibyśmy $h(rx)=rh(x)$.
Tymczasem $h(E_{21})=E_{12}E_{21}=E_{11}$,
natomiast $E_{21}h(1_R)=E_{21}E_{12}=E_{22}$.
Ponieważ $E_{11}\ne E_{22}$, $h$ nie jest
$R$-liniowe. Wyjaśnia to, dlaczego w
Stwierdzeniu~\ref{aa:st:modules-hom-abelian}
przemienność pierścienia jest założeniem, a nie
jedynie wygodną konwencją.
\end{przyklad}

\section{Dlaczego strony modułu mają znaczenie?}
\begin{stwierdzenie}[Endomorfizmy lewego modułu regularnego]
\label{aa:st:modules-end-regular}
Dla pierścienia $R$ z jedynką zachodzi izomorfizm
pierścieni $\operatorname{End}_R({}_RR)\cong R^{\mathrm{op}}$,
gdzie ${}_RR$ oznacza lewy $R$-moduł dany przez
mnożenie w pierścieniu, a w pierścieniu endomorfizmów
mnożeniem jest składanie odwzorowań.
\end{stwierdzenie}
\begin{proof}
Jeśli $f:R\to R$ jest lewym $R$-liniowym endomorfizmem,
to dla każdego $r\in R$ mamy
$f(r)=f(r\cdot1_R)=rf(1_R)$. Zatem $f$
jest jednoznacznie wyznaczony przez element
$a=f(1_R)$ i ma postać $f_a(r)=ra$.
Każdy taki wzór określa lewy $R$-liniowy endomorfizm,
bo $f_a(sr)=(sr)a=s(ra)=sf_a(r)$.
Dla $a,b\in R$ mamy
$(f_a\circ f_b)(r)=f_a(rb)=(rb)a=r(ba)=f_{ba}(r)$.
Zatem odwzorowanie $a\mapsto f_a$ zachowuje
mnożenie w pierścieniu przeciwnym $R^{\mathrm{op}}$,
a także dodawanie i jedynkę. Jest bijekcją,
więc jest izomorfizmem pierścieni.
\end{proof}

\begin{uwaga}[Dalszy program teorii modułów]
\label{aa:uw:modules-future}
Zdefiniowaliśmy moduły, podmoduły, prezentacje,
krótkie ciągi dokładne i moduły rzutowe, nie zakładając
przemienności pierścienia. W dalszym ciągu kursu
opracujemy klasyfikację modułów nad dziedzinami
ideałów głównych oraz postać Smitha. Osobno
rozwiniemy \emph{iloczyn tensorowy modułów}, który
nie jest ani sumą prostą, ani iloczynem prostym.
Dla nieprzemiennego pierścienia iloczyn tensorowy
nad $R$ wymaga co do zasady prawego $R$-modułu
i lewego $R$-modułu; przyjmiemy wtedy odpowiednią
własność zbilansowania.
\end{uwaga}

\setlength{\parskip}{4pt}
